% SPDX-FileCopyrightText: 2026 Bart van de Lint
% SPDX-License-Identifier: AGPL-3.0-or-later
\section{Method}
\label{sec:model}

The model turns a ship design into a cost per tonne-kilometre on a route, the full cost of carrying one tonne of cargo one kilometre. A design is a hull of a given size, a drive and at most one rig. The drive is a mechanical diesel engine on very low sulphur fuel oil (VLSFO), today's fuel, or a mechanical dual-fuel engine on e-ammonia, the zero-emission fuel of this paper (\cref{sec:fuel_choice}). The rig is a towing kite, a set of rigid wing sails or a set of Flettner rotors. \Cref{fig:pipeline} shows the chain. The ship models give the thrust the propeller must deliver in every wind, wave and speed. A table of these drives a weather router, which chooses each ship's path and speed through three years of recorded weather. The router's passages give fuel, time and kite flight hours, and the cost model turns those into a cost per tonne-kilometre. The cost model also sets the hourly price of time that the router trades against fuel. Every input is listed in one table with a value, a low and a high bound, a unit and its source, or the statement that it has none (\cref{sec:inputs}). The code is published open source and is the full record of every number and derivation.\tbd{ Repository link.}

This section says what each step does, what it assumes and why. \Cref{sec:details} gives the formulas, calibrations and sources in the same order.

\begin{figure*}[tbp]
\centering
\begin{tikzpicture}[
  node distance=4mm and 4.5mm,
  box/.style={draw, rounded corners=2pt, align=center, font=\footnotesize, minimum height=11mm, text width=21mm, fill=#1},
  box/.default=white,
  data/.style={draw, align=center, font=\footnotesize, minimum height=11mm, text width=20mm, fill=gray!8},
  arrow/.style={-{Stealth[length=2mm]}, thick}]
\node[data] (design) {Design\\ length, drive, rig};
\node[data, below=of design] (inputs) {Inputs\\ value, range, source};
\node[box=blue!8, right=of design] (ship) {Ship models\\ hull, waves, propeller, engines};
\node[box=orange!12, below=of ship] (rig) {Rig force tables\\ kite eights, wings, rotors};
\node[box=blue!8, right=of ship, yshift=-7.5mm] (balance) {Force balance\\ leeway, rudder, heel; best trim};
\node[box=blue!8, right=of balance] (table) {Performance table\\ thrust, fuel power};
\node[data, below=of table] (weather) {ERA5 wind and waves, 2017--2019};
\node[box=green!10, right=of table] (router) {Weather router\\ dynamic programme};
\node[data, below=of router] (corridor) {Route corridors, land mask};
\node[box=red!8, right=of router] (cost) {Cost model\\ fuel, time, ports $\rightarrow$ USD/tkm};
\draw[arrow] (design) -- (ship);
\draw[arrow] (design.south east) -- (rig.west);
\draw[arrow] (ship) -- (balance);
\draw[arrow] (rig) -- (balance);
\draw[arrow] (balance) -- (table);
\draw[arrow] (table) -- (router);
\draw[arrow] (weather) -- (router);
\draw[arrow] (corridor) -- (router);
\draw[arrow] (router) -- (cost);
\draw[arrow] (inputs) -- (rig);
\end{tikzpicture}

\caption{The model, step by step, each box with its section. Arrows carry what one step hands the next, and dashed frames mark what is repeated, as labelled. Blue boxes model the ship, orange the rig, green the voyage and red the cost; grey boxes are data.}
\label{fig:pipeline}
\end{figure*}

\subsection{Ships}
\label{sec:ships}

We use two sets of ships. One \SI{85}{m} coaster scaled geometrically to \SI{200}{m} isolates the effect of size, spanning \num{3800} to \num{50000}~t of deadweight, the mass of cargo, fuel and stores a ship can carry. Thirteen standard classes of merchant ship (\cref{tab:classes}) run from the same coaster to a container ship of \num{24000}~TEU, that is \num{24000} twenty-foot containers. Their dimensions, deadweight, speed and engine power come from MAN Energy Solutions' propulsion-trend papers and named ships, their prices from broker reports and contracts of 2025 and 2026, and their cargo values from commodity prices. Each hull floats at its design draught, the loaded draught its engine and propeller are designed for. Its deadweight is that at the deeper scantling draught its structure allows, since that sets the cargo it can lift. Each class keeps the engine power of its source design, and the router chooses the speed.

\begin{figure*}[tbp]
\centering
\plot[\textwidth]{ship_classes}
\caption{The thirteen standard ship classes in deadweight order, seen from the side to one scale, each with the best kite, wing sails and rotors on VLSFO, the geometric mean over the six routes of the best area. Hull, deck height, deckhouse and the rigs are drawn as the model places them: deck rigs spread over the deck below the air draught, a single unit at the hull's lateral centre and the kite's launch mast at the forward end of the deck. The kite is drawn to scale across its span just above its mast; its far longer tether is not.}
\label{fig:ship_classes}
\end{figure*}

The classes fall into two groups that pay for a rig in different ways (\cref{tab:classes}, column Limit). A deadweight-limited ship fills its deadweight before its space. Every tonne of rig costs it a tonne of cargo, while the deck the rig stands on carries no cargo. The coaster, the four bulk carriers, the three tankers and, narrowly, the Megamax container ship are deadweight-limited. A space-limited ship fills its slots or car decks first. The mass of its rig costs no cargo, but every slot the rig blocks is cargo lost. The feeder, Panamax and Neo-Panamax container ships are space-limited, and so is the car carrier, whose rig stands on a weather deck that carries no cars and so costs it nothing at all. \Cref{sec:deck_rigs} gives the model of both limits and where each class lies.

\begin{table*}[tbp]
\centering
\footnotesize
\caption{The thirteen standard ship classes, ordered by deadweight as in all class figures: waterline length $L$, beam $B$, design draught $T$, block coefficient $C_B$ at that draught (the underwater volume over that of a box of the same length, beam and draught), deadweight at the scantling draught, slot or deck capacity of the space-limited ships in twenty-foot containers (TEU) or standard cars (CEU), the limit the cargo of the class fills first, service speed and main-engine power of the source design, the factor $f_R$ that calibrates the calm-water resistance to it, newbuild price, cargo value, container tiers on deck, the lowest air draft (height above the waterline) on common calls, the share of calls a fixed rig above it loses and the time a fixed rig costs cargo handling per call ($\infty$ where it prevents it), and the frontal and lateral areas above the waterline with the deck cargo; \cref{sec:deck_rigs} explains the air draft, the lost calls and the delay. Where a source gives only the scantling draught, as for the container ships and the Aframax tanker, the design draught is the scantling draught times the ratio of the two in a published design of the same type. Each value has its own source.}
\label{tab:classes}
\setlength{\tabcolsep}{3pt}
\begin{adjustbox}{max width=\textwidth}
\begin{tabular}{@{}lrrrrrrrrrrrrrrr@{}}
\toprule
Class & $L$ & $B$ & $T$ & $C_B$ & Deadweight & Capacity & Speed & Power & $f_R$ & Price & Cargo & Deck & Air draft & Lost & Fixed-rig \\
 & [m] & [m] & [m] & & [t] & & [kn] & [MW] & & [MUSD] & [USD/t] & tiers & limit [m] & calls [\%] & delay [h] \\
\midrule
general cargo coaster & \num{85} & \num{13.2} & \num{5.5} & \num{0.83} & \num{3800} & -- & \num{10.5} & \num{1.5} & \num{1} & \num{11.6} & \num{2500} & \num{0} & \num{35} & \num{10} & $\infty$ \\
feeder container ship & \num{165} & \num{25.6} & \num{8.75} & \num{0.66} & \num{21000} & \num{1600} TEU & \num{18.9} & \num{10.5} & \num{0.86} & \num{33} & \num{5120} & \num{4} & \num{40} & \num{5} & $\infty$ \\
PCTC car carrier & \num{192} & \num{37.62} & \num{9.35} & \num{0.62} & \num{25563} & \num{9100} CEU & \num{14.5} & \num{13.7} & \num{0.94} & \num{100} & \num{15000} & \num{0} & \num{53} & \num{5} & \num{0} \\
Handysize bulk carrier & \num{177} & \num{30.5} & \num{10.3} & \num{0.8} & \num{42000} & -- & \num{14} & \num{6.5} & \num{0.86} & \num{31} & \num{250} & \num{0} & \num{45} & \num{6} & \num{2} \\
MR product tanker & \num{174} & \num{32.2} & \num{11} & \num{0.81} & \num{47000} & -- & \num{15.1} & \num{9.3} & \num{0.87} & \num{51} & \num{650} & \num{0} & \num{45} & \num{6} & \num{0} \\
Panamax container ship & \num{245} & \num{32.2} & \num{11.2} & \num{0.65} & \num{51000} & \num{4000} TEU & \num{21.1} & \num{22.8} & \num{0.91} & \num{58} & \num{5120} & \num{6} & \num{53} & \num{5} & $\infty$ \\
Ultramax bulk carrier & \num{196} & \num{32.2} & \num{12.2} & \num{0.86} & \num{64000} & -- & \num{14.5} & \num{8.6} & \num{0.75} & \num{34.5} & \num{180} & \num{0} & \num{45} & \num{6} & \num{2} \\
Kamsarmax bulk carrier & \num{225} & \num{32.26} & \num{12.2} & \num{0.88} & \num{82000} & -- & \num{14.5} & \num{9.7} & \num{0.8} & \num{37.5} & \num{170} & \num{0} & \num{45} & \num{6} & \num{4} \\
Aframax tanker & \num{236} & \num{44} & \num{13.9} & \num{0.8} & \num{114824} & -- & \num{14.15} & \num{10.3} & \num{0.87} & \num{75} & \num{506} & \num{0} & \num{45} & \num{6} & \num{0} \\
Neo-Panamax container ship & \num{350} & \num{50.8} & \num{14.3} & \num{0.67} & \num{152000} & \num{14000} TEU & \num{22.4} & \num{49.4} & \num{0.87} & \num{152.5} & \num{5120} & \num{8} & \num{53} & \num{5} & $\infty$ \\
Capesize bulk carrier & \num{282} & \num{45} & \num{16.5} & \num{0.84} & \num{180000} & -- & \num{14.6} & \num{16.2} & \num{0.81} & \num{75} & \num{100} & \num{0} & \num{60} & \num{1} & \num{4} \\
Megamax container ship & \num{398} & \num{61.5} & \num{14.8} & \num{0.72} & \num{230000} & \num{24000} TEU & \num{21.7} & \num{60} & \num{0.85} & \num{267.3} & \num{5120} & \num{10} & \num{57} & \num{5} & $\infty$ \\
VLCC & \num{320} & \num{60} & \num{20.5} & \num{0.81} & \num{301781} & -- & \num{15.5} & \num{27.2} & \num{0.99} & \num{130} & \num{506} & \num{0} & \num{62} & \num{1} & \num{0} \\
\bottomrule
\end{tabular}
\end{adjustbox}

\end{table*}

\paragraph{Calm water} The resistance $R_\mathrm{cw}(V)$, the force needed to push the hull through calm water at speed $V$, follows Holtrop and Mennen \citep{holtrop1982,holtrop1984}\tbd{ (sources not read)}, an empirical method fitted to model tests of many hulls. Uncalibrated, it asks for 9 to 71\,\% more power at the service speed than the class's source design implies. We therefore scale the resistance of each class by one factor $f_R$ (\cref{tab:classes}). In calm water the model then needs the power the source design delivers at its service speed, less the sea margin, the extra power a designer holds back for wind, waves and fouling. Without this factor every class would burn more fuel than its real design (\cref{sec:det_ships}).

\paragraph{Wind on the ship} The ship above water, with its deckhouse and deck cargo, also meets the wind. The force along the heading is the air resistance $R_\mathrm{air}$ below, and the side force and its moments add to the rig's. Coefficients come from the wind-tunnel measurements of \citet{blendermann1993} on models of each ship type, with the bulk carriers taking the tanker's since he measured none. The areas come from one of his ships scaled to the class, with extra container tiers on deck (\cref{tab:classes}). The ship keeps its design draught and, for the container ships, its full deck cargo on every leg, so loaded and ballast passages are not told apart (\cref{fig:wind_load,sec:det_ships}).

\begin{figure*}[tbp]
\centering
\plot[\textwidth]{wind_load}
\caption{Coefficients of the wind load on the ship itself over the apparent wind angle, on its lateral area, for one class of four types with their deck cargo, from the loading functions of \citet{blendermann1993}: (left) along the ship, positive forward, and (right) to the side. The other container ships and bulk carriers follow the Megamax and the Aframax tanker.}
\label{fig:wind_load}
\end{figure*}

\paragraph{Waves} Waves add a resistance that grows with the square of the significant wave height $H_s$, the mean height of the highest third of the waves. It comes from a semi-empirical formula for regular waves \citep{mittendorf2022}, summed over the frequencies of a standard sea spectrum with the height, period and direction of the ERA5 weather record (\cref{sec:det_ships}).

\paragraph{Propeller and engines} The propeller is an ideal disk that accelerates the water through it, with a viscous loss. Its efficiency falls as the thrust per unit of disk area rises, so a rig that takes load off the propeller also makes it work more efficiently. The engine burns fuel at a peak efficiency of 44\,\% of the fuel's heating value, VLSFO and e-ammonia alike, which falls at part load. The fuel-oil and e-ammonia drives each have one engine. The e-methanol drive of \cref{sec:res_fuels} is diesel-electric, with three generator sets, and runs the number that burns least. Every drive gives the propeller what the main engine of the class's source design gives at its rating, plus \SI{60}{kW} for the ship's own needs. Speeds that need more are infeasible (\cref{sec:det_ships}).

\paragraph{The zero-emission fuel}
\label{sec:fuel_choice}
We compare the rigs on one zero-emission fuel for every ship, the one expected to carry most of deep-sea shipping by 2050, so that the comparison isolates what wind is worth to a ship on e-fuel. Of e-ammonia, e-methanol and e-LNG, the net-zero pathways name ammonia. IRENA's 1.5\,\textdegree{}C pathway gives it 43\,\% of shipping's energy in 2050 and calls it the backbone of the sector's decarbonisation \citep{irena2021ship}. The Mærsk Mc-Kinney Møller Center's pathway gives it just over half and finds it the cheapest e-fuel to make, because it needs no carbon \citep{mmmczcs2021its,mmmczcs2021fuel}. Today's orders point the other way. LNG holds 36.8\,\% of the tonnage on order, methanol 10.4\,\% and ammonia 0.7\,\%, and DNV, whose orderbook that is, sees no silver bullet among the fuels \citep{dnv2025mf}. We read the LNG orders as a bet on a transition fuel, since no net-zero pathway we read makes e-LNG a main fuel of 2050. Ammonia engines are the least mature of the three, and we assume engines of the power of the largest container ships will exist when the fuel does. \Cref{sec:res_fuels} shows how much the choice of e-fuel moves the savings and \tbd{that it leaves the ranking of the rigs unchanged}.

The e-ammonia engine burns 5\,\% of its energy as VLSFO pilot fuel, priced and counted for CO$_2$ as VLSFO. Its tank takes 3.6 times the volume of a fuel-oil tank of the same energy and, filled, weighs about twice its fuel, our estimate. Every fuel store is sized for the longest leg in the weather sample with a 30\,\% reserve, so a rig that saves fuel on that leg also shrinks the tank. On the container ships and the car carrier the tank takes cargo space (\cref{sec:deck_rigs,sec:det_ships}).

\subsection{Balance of rig forces}
\label{sec:balance}

A rig does not only pull forward. Its force has a component along the ship's heading, the thrust $F_x$, and one across it, the side force $F_y$. It also leans the ship over with a heeling moment $M_x$ about the waterline and turns it with a yaw moment $M_z$ about the middle of the hull, and a rotor draws a power $P_\mathrm{rig}$ to spin. The wind on the ship adds its own side force and moments, so in what follows $F_y$, $M_x$ and $M_z$ are the rig's plus the ship's. The hull answers the side force by moving slightly sideways, at a leeway angle $\beta$ of a few degrees, so that it acts as a low, long wing. Like any wing it then has an induced drag $R_\beta$. The rudder answers the yaw moment and adds its own drag $R_r$. The thrust the propeller must give is
\begin{equation}
T_\mathrm{req} = R_\mathrm{cw} + R_\mathrm{air} + R_\mathrm{aw} + R_\beta + R_r - F_x,
\label{eq:thrust}
\end{equation}
the calm-water, air and wave resistances, plus the costs of side force, less the rig's thrust. A rig setting is not used if it heels the ship more than \SI{8}{\degree} \citep{tillig2020}, needs a leeway above \SI{10}{\degree} or a rudder lift coefficient above 1. Where the wind on the ship alone already passes one of these limits, a setting may not go further past it. Of the remaining settings the ship uses the one that needs the least engine output, so a rotor spins faster only when its extra thrust saves more propeller power than its spin costs. Each trim, the wing's lift, the rotor's spin and the kite's weight on side force below, is solved for as a continuous quantity rather than picked from a list, so every rig sits equally close to its own optimum (\cref{sec:det_balance}).

Where the rig stands matters. A single wing or rotor stands at the hull's lateral centre, the point where the hull's own side force acts, so its side force costs only leeway and heel. The soft kite tows from its launch mast at the forward end of the deck, as the ship kites built so far do \citep{fritz2013} and kite studies assume \citep{traut2014}, so the rudder holds a small turning moment. \Cref{sec:tow_point} compares this with tow points further aft. The VTOL kite tows from the lateral centre. A row of several wings or rotors is centred at the middle of the deck, and the rudder carries the yaw moment that remains. The kite pulls from \SI{10}{m} above the waterline on the reference ship and the deck rigs at their mid height, so for the same side force the kite heels the ship less. \Cref{sec:det_balance} gives the hull and rudder model, a regression fitted to model tests \citep{kim2021} and a rudder sized by classification rules \citep{dnv_rudder}. \Cref{fig:force_balance} draws the balance for each rig on one ship in a beam wind.

\begin{figure*}[tbp]
\centering
\plot[\textwidth]{force_balance}
\caption{Horizontal forces on the \SI{140}{m} ship at \SI{10.5}{kn} in a \SI{10}{m/s} true wind on the beam, seen from above with the bow up, to scale, for each rig at its best trim. The rig force acts at the kite's towing point or at the centre of the wings or rotors, and the wind load on the ship above the waterline at its own centre; the hull answers their side force at its lateral centre, and the rudder carries what yaw moment remains. The drag is the calm-water and wave resistance with the drag of leeway and rudder, and the propeller gives what the rig and the wind load leave of it (\cref{eq:thrust}). The dashed line is the kite's tether.}
\label{fig:force_balance}
\end{figure*}

\subsection{Wing sails and rotors}
\label{sec:deck_rigs}

Wings and rotors stand in a row on the centreline. A rig is split into as few units as the \SI{45}{m} air draft above the waterline and the tallest units built allow, spread evenly over 75\,\% of the deck length. Units close together shade each other, so a larger rig gains less than its area. We fit this loss to wind-tunnel tests and flow simulations of two to four rotors and three and four wing sails \citep{bordogna2020,han2025,malmek2024,xu2025}, and take back a third of the wings' loss for sheeting each sail on its own \citep{vanreen2026}. Both rigs take the wind at their mid height, from a wind profile fitted to the weather record over the study region (\cref{sec:disc_profile}).

A wing sail is a rigid upright wing with an aspect ratio, span squared over area, of 2, trimmed like a sail by its lift coefficient $C_L$ between 0 and 2.0, or stalled at a steep angle to the wind running before it. A Flettner rotor is a tall spinning cylinder with end discs. The spin turns the passing wind into a large side force, the Magnus effect, which the rotor points forward. Its lift, drag and spin power follow the measured curves of the ShipCLEAN model over the spin ratio, the rotor's surface speed over the wind speed \citep{tillig2020}. At a spin ratio of 3 its lift coefficient is about five times what the wing reaches. The rotor is trimmed by its spin and stopped when the apparent wind exceeds \SI{20}{m/s}. A rotor that does not spin still has drag, so a rotor ship needs slightly more power in calm water (\cref{sec:det_deck}).

\paragraph{Rig mass and deck space} A larger rig costs cargo in two ways. Its mass comes off the deadweight, and the deck rigs need deck length, with a height the air draft limits, so a larger rig means more units rather than taller ones. Whatever does not fit is not a design. On container ships each deck unit also blocks the deck stacks its footprint covers over the length of one container. We estimate this from the geometry, since no container ship carries a wing or rotor yet. We size every rig by searching its area for each ship class, drive and route, so each is compared at close to its own best size (\cref{sec:det_deck}).

\paragraph{Mass or volume limit} A ship carries cargo up to its deadweight or up to its space, whichever comes first (\cref{sec:ships}). On a deadweight-limited ship every tonne of rig or fuel store costs a tonne of cargo. On a space-limited ship a rig costs cargo only where it blocks a slot, or where its mass and the fuel store take the ship to its deadweight. Mean loads per container and per car come from European port statistics \citep{eurostat_mar_mg_am_cvh,eurostat_maritime_manual}. The line between the groups is close for the Megamax and the Neo-Panamax (\cref{sec:det_deck}), and \cref{sec:res_groups} gives what moving a class across it does to the result.

\paragraph{Fixed or retractable} Wings and rotors either stand fixed or can be tilted, folded or moved on rails, which costs 12\,\% more and takes minutes in port. We give each ship whichever costs less. A fixed rig costs time in port at the ship's own cost per hour, because it obstructs cargo handling, from a few hours per call on a bulk carrier to prevented handling on a container ship. Its height also keeps an empty ship, which floats higher, from some berths behind low bridges. These delays and shares are our estimates, as no operator publishes them (\cref{tab:classes,sec:det_deck}). The kite's launch mast is always telescopic, as on every ship kite built.

\subsection{Kite}
\label{sec:kite}

\paragraph{Shape and aerodynamics} The kite is a ram-air wing like a paraglider: two skins of cloth with ribs between them, inflated through openings at the front, curved into an arc in flight (\cref{fig:kite}). The base kite has a projected aspect ratio of 3, and every kite area in this paper is projected. Its lift and drag come from a vortex model of the whole wing \citep{cayon2023}\tbd{ (source not read)} on section data from a neural-network airfoil code \citep{sharpe2025} (\cref{fig:polars}), for the design shape. A kite in flight does worse than its design shape. Its glide ratio, lift over drag, sets how fast it can fly across the wind and so how hard it pulls. A rigid scale model reaches 8 to 8.5 in a wind tunnel \citep{poland2025}, but about 4 in flight with its lines and control unit \citep{oehler2019}, and the SkySails ship kites reached 3.9 to 5 with their tether \citep{erhard2013b,paulig2013}. We add the difference as one drag coefficient of 0.08, which brings the base kite to the glide ratio of flown kites. A kite that keeps its design shape is the upper bound. We hold the aspect ratio at 3, a value used by industry for which we have no public reference. A more slender kite could fly more efficiently but needs larger launch gear and is less reliable, which the model does not price, and the glide ratio is calibrated to flown kites rather than taken from the shape (\cref{sec:det_kite}). We report the span with every kite.

\paragraph{Lines} A net of bridle lines spreads the tether force over the canopy and meets at a control unit, a pod that steers the kite by pulling lines. The single tether to the ship is 13 spans long, as on the larger SkySails kite \citep{paulig2013}, and its length is held during flight and changed only at launch and landing. Its rope is sized for the force limit below with a design factor of 3 on its breaking load, and the point where it is made fast on the ship turns with it, so the rope is not bent back and forth as the kite flies. No ship or kite tether life has been measured. We estimate it from creep, the slow stretching of a rope held under load, counting the load peaks of gusts and of the bow moving in waves. This gives \num{16000} flight hours, with a range of \num{4500} to \num{45000} (\cref{sec:det_kite}).

\paragraph{Glide ratio of the whole system} What pulls the ship is the glide ratio of the kite with its lines. The drag of the bridle, of the part of the tether that acts at the kite and of flight brings it to $G=3.9$ at every size, within the 3.5 to 4.6 that published ship-kite studies take \citep{traut2014,cedelft2016,podeur2016}. A kite cannot hold a steady force further upwind than the edge of its wind window, so it gives forward thrust only where the apparent wind comes from at least \SI{14}{\degree} to \SI{16}{\degree} off the bow, and little close to that limit (\cref{sec:det_kite}).

\paragraph{Canopy} The canopy is sewn from polyester ripstop of \SI{200}{g/m^2} and weighs \SI{0.75}{kg} per projected square metre with its seams and reinforcements. A kite built for a higher force limit gets heavier cloth, seams and reinforcements (\cref{sec:det_kite}). Its life is a goal of \num{3000} flight hours (\cref{sec:canopy_life}). Its price is a goal of USD~150 per projected square metre, against USD~234 today, built up from the retail price of cloth and the sewing time of paragliders (\cref{fig:canopy_price}). The goal assumes cloth bought by the roll from the mill and the sewing time of a large series of military cargo parachutes. The step from retail to mill price is our judgement and the weakest part of the goal. Wings, rotors and tethers have no such goal, since they are built up from industrial costs or bulk prices rather than retail ones.

\begin{figure*}[tbp]
\centering
\plot[\textwidth]{canopy_price}
\caption{Canopy price per projected square metre at a cloth weight of \SI{200}{g/m^2}. (a) Today's price, the goal with cloth still bought at retail, and the goal, each split into cloth, sewing with lines and hardware, and the margin of the market factor. (b) Price of a sewn square metre of cloth against cloth weight, at retail today and from the mill at the goal; the vertical line marks the canopy's cloth weight.}
\label{fig:canopy_price}
\end{figure*}

\paragraph{Flight} The kite is steered by the heading of its flight through the air, as the SkySails autopilot flies it \citep[Sec.~VIII]{erhard2013}. At each instant the model finds the air speed at which lift, drag, weight, tether pull and steering force balance. This quasi-steady model neglects the kite's momentum from one instant to the next. Without mass and flying straight across the wind it reproduces the classic result of \citet{loyd1980}\tbd{ (source not read)} to $10^{-6}$ (\cref{sec:checks}). Where the tether force would exceed the structural limit of \SI{0.5}{kN} per square metre of kite, the kite depowers, lowering its angle of attack until the force fits.

The kite holds a heading until it reaches one end of its pattern, turns and flies back, and the force on the ship is the tether force averaged over one cycle. Seen from the ship, the kite flies in the wind window, the quarter sphere downwind of the tether's anchor (\cref{fig:kite_frames}). At its edge the kite hangs nearly at rest, and in its centre it flies fastest and pulls hardest. One family of patterns covers slow sweeps along the edge, fast figure eights across the wind and rest at the edge (\cref{fig:patterns,sec:det_kite}).

\paragraph{Pattern choice} For each apparent wind speed and ship course a search picks the pattern, its position and width in the window, its heading and the kite's angle of attack that give the most thrust along the course, less a weighted penalty on the side force. The weight, any value from 0 to 3, is the kite's trim. A larger weight gives less thrust and less side force, which the ship chooses where heel or the rudder limit it. Downwind the best patterns are figure eights, and towards the wind they are sweeps along the edge of the window (\cref{fig:patterns}). The search is local and does not always find the best pattern. It scatters one ship's cost saving by about 0.34 points (one standard deviation), so smaller differences between kites on one ship are not resolved (\cref{sec:det_kite}).

\paragraph{Limits} A kite flies from the lowest wind of \cref{sec:min_wind} up to an apparent wind of \SI{25}{m/s}, the limit of current kite systems \tbd{cite}. Its launch is limited to \SI{17}{m/s}, its landing to \SI{20}{m/s}, and launch, landing and flight to waves of 3, 4 and \SI{5}{m}. The landing wind is the design limit of the stowage of a full-size land kite \citep{skypower100_tib}. No source gives the launch wind or the wave limits, so they are our estimates with wide ranges. A kite may always come down, but a landing outside its limits is forced, with an estimated 5\,\% chance of losing the kite, so the router brings it down ahead of a storm where it can. Every leg ends with the kite down. In the base case a kite also launches at night, and a case with launches by day only tests that assumption.

\paragraph{Launch and landing} One launch and landing costs crew time, a small chance of a retry, of a tear that ends the canopy and of losing the kite in the water, the mean cost of turning the ship for a launch in a head wind, and the wear of the tether as it is reeled in and out. The chance of a tear or a loss is together about 1 in \num{10000} per cycle as a goal, against 0.2\,\% for a tear and 0.05\,\% for a loss that we would estimate today. For \SI{25}{min} of launch and \SI{30}{min} of landing and stowing the ship gets no force from the kite, on reports of 10 to \SI{20}{min} each for a kite of \SI{160}{m^2} \citep{konstruktionspraxis_skysails}. No source gives rates of retries, tears or losses per launch for a ship kite, so these are goals or estimates with wide ranges (\cref{tab:inputs,sec:det_kite}).

\begin{figure*}[tbp]
\centering
\plot[\textwidth]{kite_frames}
\caption{Frames and angles of the kite model. (a) The wind window seen from the tow point looking downwind: the kite at roll angle $\phi$ and angle $\vartheta$ off downwind, the unit vectors $\hat{\mathbf e}_\vartheta$ towards the edge of the window and $\hat{\mathbf e}_\phi$ along the roll, and the heading $\hat{\mathbf h}$ at the angle $\psi$ between them. The shaded wedge is a band of roll angles $\phi_0\pm\phi_a$ a pattern sweeps; no part of a pattern flies in the grey band below the lowest elevation. (b) The same seen from the side. (c) Seen from above, the mean tether force $\overline{\mathbf F}$ and its components along and across a course $\chi$ off the apparent downwind direction; the kite pulls the ship forward wherever the force leans ahead of the beam of that course.}
\label{fig:kite_frames}
\end{figure*}

\begin{figure*}[tbp]
\centering
\plot[\textwidth]{kite_geometry}
\caption{Planforms of the kites of projected aspect ratio 3, 5 and 8 seen from above, to scale, at a projected area of \SI{100}{m^2}, with their projected span. The leading edge is at the top.}
\label{fig:kite}
\end{figure*}

\begin{figure*}[tbp]
\centering
\plot[\textwidth]{polars}
\caption{Lift, drag and glide ratio of the kites of projected aspect ratio 3, 5 and 8 in their design shape over angle of attack, from the Vortex Step Method, without bridle, tether and flight drag.}
\label{fig:polars}
\end{figure*}

\begin{figure*}[tbp]
\centering
\plot[\textwidth]{patterns}
\caption{Optimal flight paths of the base \SI{800}{m^2} kite over one cycle, as elevation over azimuth off the apparent downwind direction, coloured by the ship's course off the apparent downwind direction, in an apparent wind of \SI{6}{m/s} (a) and \SI{12}{m/s} (b), without side-force weight, and the time-averaged thrust and side force they put on the ship (c, d). Downwind the kite flies figure eights through the wind window. Towards the wind it sweeps along the edge of the window, and where a kite at rest would serve the ship better no path is drawn.}
\label{fig:patterns}
\end{figure*}

\subsection{Lowest wind for flight and launch}
\label{sec:min_wind}

A kite needs enough wind to carry its own weight, so the lowest wind decides how often a kite can fly on a calm route. Three thresholds follow from a balance of forces in steady flight, derived in \cref{sec:det_min_wind} beside their nearest published relatives \citep{bauer2018,fagiano2017takeoff,joshi2024powercurve}: $V_\mathrm{cross}$, the lowest wind in which a kite crossing the wind stays up, $V_\mathrm{climb}$, the lowest wind in which it can also climb, and $V_\mathrm{lift}$, the wind a kite at rest needs to lift off. A launch also needs at least \SI{5}{m/s}, measured on a full-size land kite \citep{skypower100_tib}, and Beaufort~3, the wind given for the SkySails ships in perfect conditions \citep{paulig2013,noaa2009}. \Cref{fig:min_wind} gives these thresholds over kite area. The launch wind is the reported \SI{5}{m/s} up to the largest kites, where the climb takes over. The first flyable pattern appears at \tbd{5 to \SI{6.5}{m/s}} on the best course and at \tbd{8 to \SI{9}{m/s}} upwind, so thrust starts no lower than the launch wind. We do not let the ship steam faster to launch a kite, because no pattern gives thrust below that wind (\cref{sec:det_min_wind}).

\begin{figure}[tbp]
\centering
\plot{min_wind}
\caption{Lowest apparent wind against projected area for the kite of aspect ratio 3 at the lowest elevation: the crosswind bound $V_\mathrm{cross}$ (\cref{eq:min_cross}), the climb $V_\mathrm{climb}$ (\cref{eq:min_climb}), lift-off with its margin $V_\mathrm{lift}$ (\cref{eq:lift_off}), and the launch wind the model uses, the largest of the last two and the reported launch wind.}
\label{fig:min_wind}
\end{figure}

\subsection{Canopy life}
\label{sec:canopy_life}

The canopy is the part of a kite that wears out in flight, and its life sets what the kite pays for each hour it flies. The deck rigs have no such running cost. We set the life as a goal of \num{3000} flight hours at \SI{200}{g/m^2}, what development can reach and not where things stand today. The model pays one canopy per \num{3000} hours flown. The range of the life is run as a sensitivity, and \cref{fig:missed_goals} gives the cost saving over the whole range, so a reader can take the result at any life.

No ship kite has a published canopy life \citep{emsa2023}, and no source gives the hours of a polyester sail or canopy together with its cloth weight and a measured load. We tried to derive a life from two things. One is the loss of cloth strength in sunlight, fitted to tests of polyester outdoors and under lamps. The other is mechanical wear from flexing, flutter, folding and chafe. No test measures wear on canopy cloth, so we take its rate from what makers of sails, balloons, paragliders and kites state of the life of their products (\cref{tab:canopy_life}). All but one of these statements are claims with no test shown. Run until the cloth keeps five times the tension it carries at the force limit, the two give \num{3800} flight hours. Wear is three quarters of that decay, so the estimate rests mainly on the makers' statements. Changing one input at a time moves the life between \num{1200} and \num{6900} hours, and we cannot narrow that range. The longest statements for any cloth, \numrange{3500}{4000} hours for a sail \citep{flynn2015} and the same for one kite \citep{kitepower2024}, are not measured either. With no measured life the goal cannot be read off the evidence. We take \num{3000} hours, on the careful side of what makers state, so that no result rests on the longest lives anyone states. It lies below those statements and below the derived life. It lies above the hours a sail keeps its shape and the life the decay gives if a canopy is retired as early as aircraft fabric and balloon cloth are (\cref{app:canopy_life_range}). A canopy that keeps its flying shape and its seams past those lives is what the goal asks of development. We keep the cloth at \SI{200}{g/m^2} rather than the lighter cloth of about \SI{160}{g/m^2} that the decay would give \num{3000} hours, so that the gap between goal and estimate is a margin on the life and not a lighter, cheaper canopy.

No source settles whether a canopy reaches the goal. We argue that it can because a ship kite pays for its canopy by the hour flown, so every hour of life goes straight to its transport cost, while a paraglider is built as light as a pilot will carry and a sail is replaced for its shape. This is our argument, not a measurement. New cloth holds 136 times the peak tension at the force limit, so strength is unlikely to end the canopy first. Leaks, seams and flutter stay open. A flutter test of polyester ripstop at several cloth weights, and the flutter a flying kite actually sees, would narrow the life more than any further reading. \Cref{sec:det_canopy_life} gives the decay model and \cref{sec:cloth_check} the tension check.

\subsection{VTOL kite}
\label{sec:rigid_kite}

The fourth rig is the VTOL kite, for vertical take-off and landing: one or several rigid wings that carry their own propellers to rise from the deck and land on it like a helicopter, on a tether that carries electric power. It uses the force balance, the pattern search, the router and the cost model of the soft kite, and differs in its lift and drag, its pattern, its mass and its prices. It only pulls, so its propellers work only in launch and landing and windmill in flight. Every input is anchored on the MX2 design of Makani, a \SI{54}{m^2} wing of \SI{26}{m} span and \SI{1852}{kg} on a \SI{300}{m} tether, which was designed and never built \citep{echeverri2020part1}. The largest wing flown with propellers on board is the \SI{32.9}{m^2} M600, for 22~h in all.

It hovers on its propellers because no other launch of a rigid wing fits a deck and starts from rest. This is also the only launch that has flown from a floating base, once, and that flight lost the kite \citep{echeverri2020part1}. Its mass grows with span to the power 2.75 (\cref{eq:rigid_mass,fig:rigid_mass}) and, unlike the canopy, not with the force limit, since no source gives how a wing's structure grows with its design load. This favours the VTOL kite at the high bound of that limit. A wing this heavy cannot fly figure eights, so it flies circles on a cone, several kites sharing one circle. No VTOL kite has a public price. Its structure, drive and tether are priced from cost models for power-generating kites at a series of 50 or more \citep{joshi2025,bauer2018}. Its life of 15 years and its loss of one kite in \num{10000} launches and landings are goals, against an M600 that lost all four of its crosswind kites in twenty crosswind flights (\cref{sec:det_rigid}).

\begin{figure}[tbp]
\centering
\plot{rigid_mass}
\caption{Mass of a VTOL kite against its wing area: the mass law of \cref{eq:rigid_mass} through the MX2 design at the base exponent and its bounds, and the published wings with propellers, Makani's Wing 7 that flew \citep{vanderlind2013} and the M600 that was built \citep{echeverri2020part1}, and without, the designs of \citet{joshi2025} and \citet{eijkelhof2020reference}. The wings differ in aspect ratio, which the law holds fixed.}
\label{fig:rigid_mass}
\end{figure}

\subsection{Performance table}

The router needs the ship's fuel power millions of times, so the model tabulates it. The table runs over speed through the water, true wind speed at \SI{10}{m} height, true wind angle and the waves' added resistance (\cref{app:perf_grid}). It stores the thrust of \cref{eq:thrust}, once with the rig at its best setting and once with the rig stowed, and the rig power. The setting is solved against the added resistance as well, so a rig with thrust to spare puts it against the waves, and beyond the table's largest added resistance the propeller makes up the rest. The added resistance comes from a separate table, and we assume that waves do not change the rig forces. Whether a kite is down or flying is left to the router. Wings and rotors cannot be stowed and stay at their best setting.

\subsection{Weather and routes}
\label{sec:routes}

Wind comes from ERA5 \citep{hersbach2020}, a reanalysis that reconstructs past weather on a global grid from observations and a weather model, as WeatherBench~2 keeps it \citep{rasp2024}, regridded to \SI{1.5}{\degree} and six-hourly, the grid on which WeatherBench~2 evaluates forecasts. Waves come from ERA5 at \SI{0.5}{\degree} and twelve hours. The years are 2017 to 2019, over the Atlantic. The land mask is the set of ERA5 wave cells without data, plus disks for the Atlantic islands this grid does not resolve (\cref{app:perf_grid}).

Six routes from Rotterdam, Europe's largest port, stand for six trades and their wind climates (\cref{fig:routes}; \cref{fig:wind_angle}a gives the wind directions met on each): to Pointe-à-Pitre in Guadeloupe through the trade winds, to New York across the North Atlantic against the westerlies, to Bilbao across the Bay of Biscay, a short-sea route with mixed winds, to Dakar along the coast of north-west Africa, to Houston in the US Gulf, Rotterdam's largest partner across the Atlantic by cargo weight \citep{eurostat_mar_go_am_nl}, and to the iron-ore port of Tubarão in Brazil across the equator, the typical voyage of a Capesize bulk carrier. Together they span voyages from \SI{1400}{km} to \SI{9200}{km} and every relative wind direction. The trades to Asia are left out, because they lead through the Suez Canal or round the Cape, which a router confined to the Atlantic cannot represent. Each return voyage carries 30\,\% of the outbound cargo. Every result averages the same 12 departures, twelve weeks apart over the three years, each sailing outbound, two days in port, and home, and the screening of the inputs uses every second of them (\cref{tab:numerics}). Rig and bare ship sail the same departures, which leaves an error of 0.4 to 1.0 points in the cost savings against all 72 departures of the sample.

The router chooses among points on a corridor, a grid laid along a reference track. The results use stages \SI{360}{nmi} apart and points \SI{240}{nmi} apart, and \cref{sec:routing} gives the error of that spacing. The corridor is wide enough that land and the edge of the weather data bound the route rather than the corridor. The coarse wind grid smears the coast and so misses the weaker wind near land, so the router doubles the cost of steps that end within \SI{50}{nmi} of land. This only steers the route, and the penalty is the same for every ship. The router does not know shipping lanes, traffic separation schemes or emission control areas, which affect ships with and without a rig alike. To Guadeloupe the outbound corridor is centred on the trade-wind route via the Canaries and the homeward one on the great circle. The other routes use the same corridor both ways.

\begin{figure*}[tbp]
\centering
\plot[\textwidth]{routes}
\caption{Routed passages on the six routes of a \SI{140}{m} e-ammonia ship without a rig (thin grey) and with a \SI{3200}{m^2} kite (coloured by route), for a winter (a) and a summer (b) departure, over the \SI{10}{m} wind five days after departure. Dotted lines are great circles. Solid lines are outbound, dashed lines homeward, each through the corridor's nodes (dots).}
\label{fig:routes}
\end{figure*}

\subsection{Weather routing}
\label{sec:routing}

For each passage the router decides three things together: the path through the corridor, the speed on each step and, on a ship with a kite, when the kite is launched and when it is landed. It picks the passage with the lowest total cost in money: the fuel, the kite's wear and launches, and the time the passage takes, charged by the hour for the ship's capital, crew, maintenance and cargo interest (\cref{fig:router}). No arrival time is imposed, so a ship whose time is expensive sails faster. The costs that do not depend on the route, port costs and the cargo carried, are fixed for a voyage, so the cheapest passage also gives the lowest cost per tonne-kilometre. The ship without a rig is routed the same way, so a rig that makes speed cheaper saves time as well as fuel (\cref{sec:res_fuel}).

The search is dynamic programming, which avoids trying every route by one rule applied stage by stage from the departure: of all the ways to reach a point at a given arrival time and with the kite up or down, keep only the cheapest. Every ship has the same number of speeds, down to \SI{4}{kn}, at which it still holds its course in a storm (\cref{sec:det_routing}). For a step of length $\ell$ sailed at speed $V$ from time $t$, the cost is
\begin{equation}
c = \frac{\ell}{V}\left[p_f\, \dot E_f(V, U_{10}, \beta, H_s, T_p, \mu) + c_w\,\delta_\mathrm{fly} + c_t\right],
\end{equation}
where $p_f$ is the fuel price, $\dot E_f$ the fuel power from the performance table for the wind speed $U_{10}$ at \SI{10}{m} and its angle $\beta$ off the bow, the significant wave height $H_s$, peak period $T_p$ and wave heading $\mu$ along the step, $c_w$ the kite's wear cost per flight hour, $\delta_\mathrm{fly}$ one where the kite flies and zero elsewhere (plus the cost of a launch and landing at a launch), and $c_t$ the cost of time per hour. Conditions in which the engine cannot give the thrust remove the step, so the router also avoids storms the ship cannot sail through.

\begin{figure*}[tbp]
\centering
\plot[\textwidth]{router}
\caption{The weather router. (a) Its grid: stages across the corridor from departure to destination, lateral nodes at each stage, the edges one node may take to the next stage, each at any of the speeds searched, and a passage. (b) A ship with a kite carries the kite's state with each node; each state holds its arrivals in time bins (ticks), keeping the exact arrival time of the cheapest (dot), and an edge may keep the kite down, launch it, fly it or land it.}
\label{fig:router}
\end{figure*}

Dynamic programming over stages, lateral points and time is established, and so is a speed chosen per step, so neither is claimed here. \Cref{tab:routers} sets this router beside those we read in full. Two things separate it from them. First, it carries the state of a rig through the search. Of these routers, those with a rig hold it at a setting fixed before the search \citep{mannarini2024,wendling2025}, switch it by a rule in the performance model \citep{mason2023}, or trim it continuously, with no on or off \citep{dhainaut2025}. The two with a kite have no launch or landing to decide or pay for \citep{bentin2018,leskinen2024}. Second, kite, wing sails and rotors all go through one search with free speed and one cost in money. The only comparison of several rigs among these held the speed at \SI{13}{kn} and had no wing sail \citep{bentin2018}. A price on time gives one point of the trade-off between time and energy that \citet{wendling2025} compute whole.

\begin{table*}[tbp]
\centering
\footnotesize
\caption{This router and the routers it is compared with, each as its publication describes it. Speed is the ship's speed through the water or, where marked, the engine power that sets it.}
\label{tab:routers}
\begin{tabular}{@{}p{2.4cm}p{3.1cm}p{2.2cm}p{2.7cm}p{3.8cm}@{}}
\toprule
Router & Search & Speed & Objective and arrival & Rig, and how the search treats it \\
\midrule
This paper & forward dynamic programming on a corridor with time bins & chosen per step & money: fuel, kite wear and launches, time at a price; arrival free & kite, wing sails, rotors; kite up or down in the search state, launch and landing charged \\
\citet{wei2012} & forward dynamic programming on stages, lateral points and time & chosen per step & fuel at a given arrival time & none \\
\citet{mason2023} & dynamic programming over position and speed, arrival time rounded & chosen per stage, 8 to \SI{14}{kn} & fuel, arrival limit & rotors, set by a rule in the performance model \\
\citet{wendling2025} & dynamic programming over position and time & engine power chosen per arc & front of time against energy & wing sails, in speed tables made before the search \\
\citet{dhainaut2025} & gradient-based optimal control, to a local optimum & continuous & fuel, maximum travel time & rotors, spin and position trimmed continuously, no on or off \\
\citet{wang2021} & genetic algorithm seeded by shortest paths and dynamic programming & eight engine power levels per leg & fuel and closeness to an arrival time & none \\
\citet{mannarini2024} (VISIR-2) & shortest path in time-varying weather, one cost per point & engine load set per run & distance, time or CO$_2$; arrival free & sailing boat, rig fixed before the search \\
\citet{marzi2024} & isofuel, which extends each route by an equal amount of fuel & constant for the voyage & fuel & none \\
\citet{bentin2016,bentin2018} & A$^*$, a shortest-path search guided by an estimate & \SI{13}{kn} held & energy, maximum travel time & rotors (2016); kite, rotors, DynaRig (2018), as force models; no launch or landing \\
\citet{leskinen2024} & shortest path on a variable grid & set by the wind, electric motor below that & least time & kite, in a table of ship speed; no launch or landing \\
\bottomrule
\end{tabular}
\end{table*}

A kite needs this state because its flight is not a matter of one step. It can launch only in narrower wind and sea than it can stay up in (\cref{sec:kite}), a launch and a landing take about an hour in which it gives no thrust, and it must be down at the end of the passage. Whether flying pays on a step therefore depends on the steps before and after, which a table of forces per step cannot hold. Wing sails and rotors have no such history and enter through the table alone. We measure what the state and the rig's place in the search are worth on the 432 passages of two ships on three routes, by the cost the router minimises (\cref{sec:det_routing}). Against the same ship without a rig, the search lowers that cost by 12.3\,\% for the kite, 13.3\,\% for the wing sails and 24.9\,\% for the rotors. When the path is found for the ship without a rig and the rig then chooses only speed and kite along it, the reductions are 9.8, 11.9 and 21.4\,\%. A comparison on a shared path would give up a fifth of the kite's reduction, a tenth of the wing sails' and a seventh of the rotors', so routing each rig for itself is part of treating them equally. Without the state, the kite is up on any step where that pays, launched and landed at no cost in time or money and in any wind and sea, the treatment closest to the two earlier routers with a kite, which have no launch or landing. Such a search reports a reduction of 14.5\,\%, 2.2 points more than the search with the state finds. Its path, sailed by the real kite, reaches 11.7\,\%. So the state mostly corrects the kite's reported result, and adds 0.6 points through the path. The charge for a launch matters little at the base launch cost: a kite searched as if launches were free, and then charged for them, reaches 12.25\,\%, launching 5.6 times a passage instead of 3.7.

The router's accuracy rests on three errors, none of which we correct. The first is the search settings. To save time the search tries only some of the speeds in coarse time bins and skips states that look too dear, and a final pass on every speed refines the path it found (\cref{sec:det_routing}). Against a search about seven times slower (\cref{app:routing_convergence}), on two ships and three routes, these settings lower the cost saving per tonne-kilometre by 0.28 points on average and at most 0.73 for the kite, by 0.25 and at most 1.18 for the rotors, and by 0.02 and at most 0.11 for the wing sails. Single passages can do worse: on 12 of the kite's 144 passages of the measurement above, the path found for the ship without a rig came out cheaper than the search's own, on four by more than 0.5\,\% of the cost of the ship without a rig and at most by 5\,\%.

The second error is the spacing of the corridor, and it is the larger. On the same ships and routes, stages \SI{180}{nmi} and points \SI{120}{nmi} apart raise the cost saving per tonne-kilometre by 2.4 points on average for the kite, at most 4.4, by 1.7 and at most 2.6 for the rotors, and by 0.5 and at most 0.9 for the wing sails, and change the cost of the ship without a rig by at most 1.3\,\%. The savings are not converged at the spacing of the results. For the kite and the rotors each halving of the spacing adds more than the one before (\cref{sec:checks}), so these figures are only a lower bound on what a finer corridor would still add; for the wing sails the change shrinks. \citet{wendling2025} likewise find results that have not clearly stabilised at their finest stage spacing. The order of the three rigs is the same at both spacings on all six pairs of ship and route, but the spacing understates kites and rotors more than wing sails.

The third error is foresight. In the base case the router knows the weather of the whole passage, as if its forecast were perfect, which no master has. A second router therefore re-plans every 24~hours on forecasts that were actually issued, and scores every passage in ERA5 (\cref{app:forecast_router}). \tbd{On one forecast run it keeps 74\,\% of the cost saving per tonne-kilometre of the wing sails and 65\,\% of the rotors', and none of the kite's within its error. On a plan that weighs the ensemble's members one by one the three keep 80, 77 and 42\,\% (\cref{sec:res_robustness}), so perfect foresight favours the kite; from a draft run.}

\subsection{Cost per tonne-kilometre}
\label{sec:cost_model}

The cost per tonne-kilometre of a design on a route is the mean round-trip cost over the cargo work (\cref{fig:cost_tree}),
\begin{equation}
c_\mathrm{tkm} = \frac{p_f E_f + c_w h_\mathrm{fly} + C_\mathrm{port} + c_t\,(t_\mathrm{sea}+t_\mathrm{port})}{m_c\,(L_\mathrm{out} + f_\mathrm{back}\,L_\mathrm{back})},
\end{equation}
with the fuel price $p_f$ and fuel energy $E_f$, the kite's wear cost $c_w$ per flight hour and its flight hours $h_\mathrm{fly}$, the port costs $C_\mathrm{port}$ of the two calls, the cost of time $c_t$ over the sea and port time $t_\mathrm{sea}$ and $t_\mathrm{port}$, the cargo mass $m_c$, the outbound and return track lengths $L_\mathrm{out}$ and $L_\mathrm{back}$, and the backhaul fraction $f_\mathrm{back}$, the share of the outbound cargo carried back (\cref{tab:inputs}). The cargo $m_c$ is limited by mass or by volume (\cref{sec:deck_rigs}). On a ship that fills its deadweight it is the deadweight less the rig mass, \SI{50}{t} of stores and the fuel store, so the fuel of a ship on e-ammonia, with half the heating value of VLSFO, weighs twice as much and costs cargo. The time cost per hour of leg is
\begin{equation}
c_t = \frac{(a + f_m)\,(C_\mathrm{ship}+C_\mathrm{drive}) + a_\mathrm{rig}\,C_\mathrm{rig} + C_\mathrm{crew} + f_\mathrm{rig} C'_\mathrm{rig}}{8760\,\epsilon} + \frac{m\,v_c\,r}{8760},
\end{equation}
the yearly cost of the ship over the hours it trades plus the interest on its cargo. Capital $C_\mathrm{ship}$, $C_\mathrm{drive}$ and $C_\mathrm{rig}$ becomes a yearly cost through capital recovery factors $a$ and $a_\mathrm{rig}$, the constant yearly payment that repays a unit of capital with interest at a discount rate $r$ of 8\,\% over the hull life of 25 years or the rig life. $C_\mathrm{crew}$ is the yearly crew cost. Maintenance costs a share $f_m$ of 3\,\% a year of the price of the ship and drive and a share $f_\mathrm{rig}$ of 2\,\% a year of $C'_\mathrm{rig}$, the rig's price without the canopy and tether, which wear already pays for, and without the airframe of a VTOL kite, which its upkeep pays for. Charging them as well would cost \tbd{0.2} points of cost saving on the \SI{140}{m} ship. The ship trades for a share $\epsilon$ of 92\,\% of the year. The cargo $m$ on board is worth $v_c$ per tonne: the outbound cargo $m_c$ over the outbound passage and one port call, and 30\,\% of it over the return. Port dues follow the Rotterdam tariff per gross tonnage, a measure of the ship's enclosed volume, and the Guadeloupe tariff per hull volume for every route. Gross tonnage follows the 1969 Tonnage Convention, from the moulded volume up to a depth of 1.32 draughts. \Cref{sec:costs} gives all prices with their sources.

\begin{figure*}[tbp]
\centering
\begin{tikzpicture}[
  node distance=3mm and 3mm,
  term/.style={draw, rounded corners=2pt, align=center, font=\footnotesize, inner sep=2pt, minimum height=9mm, text width=#1, fill=white},
  term/.default=22mm,
  routed/.style={term=#1, fill=green!10},
  routed/.default=22mm,
  edge/.style={-{Stealth[length=1.6mm]}, thin}]
\node[term=40mm, fill=red!8] (tkm) {cost per tonne-km\\ $c_\mathrm{tkm}$};
\node[term=34mm, below left=6mm and -12mm of tkm] (trip) {cost per round trip};
\node[term=40mm, below right=6mm and -12mm of tkm] (work) {cargo work\\ $m_c\,(L_\mathrm{out}+f_\mathrm{back}L_\mathrm{back})$};
\node[routed, below left=6mm and 14mm of trip] (fuel) {fuel\\ $p_f E_f$};
\node[routed, right=of fuel] (wear) {kite wear\\ $c_w h_\mathrm{fly}$ + launches};
\node[term, right=of wear] (port) {ports\\ $C_\mathrm{port}$, two calls};
\node[routed, right=of port] (time) {time\\ $c_t\,(t_\mathrm{sea}+t_\mathrm{port})$};
\node[term=40mm, below right=6mm and -10mm of work] (cargo) {cargo $m_c$: deadweight less rig, stores and fuel, or the slots and decks left};
\node[term=30mm, below=10mm of time, xshift=-20mm] (capital) {capital\\ ship, drive, rig};
\node[term=16mm, right=of capital] (crew) {crew};
\node[term=22mm, right=of crew] (maint) {maintenance\\ $f_m C$};
\node[term=26mm, right=of maint] (interest) {cargo interest\\ $m\,v_c\,r$};
\draw[edge] (trip) -- (tkm);
\draw[edge] (work) -- (tkm);
\foreach \n in {fuel, wear, port, time} \draw[edge] (\n.north) -- (trip.south);
\draw[edge] (cargo.north) -- (work.south);
\foreach \n in {capital, crew, maint, interest} \draw[edge] (\n.north) -- (time.south);
\node[font=\scriptsize, align=left, anchor=north west] at ([yshift=-2mm]fuel.south west) {\tikz\node[draw, fill=green!10, minimum size=2mm, inner sep=0]{}; from the routed passages};
\end{tikzpicture}

\caption{How the cost per tonne-kilometre is built. Shaded terms come from the routed passages; the others from the design and the prices of \cref{sec:costs}. \Cref{fig:breakdown} gives the parts for the studied designs.}
\label{fig:cost_tree}
\end{figure*}

\textbf{Rigs out of service.} Each unit of a rig is out of service for a share $u$ of the round trips though the wind would allow it, independently of the other units. Every rig is routed in two states over the same departures and the costs of both are averaged: all units working, on the share $(1-u)^n$ of the round trips for $n$ units, and otherwise with a unit down, holding the mean share $1-u/(1-(1-u)^n)$ of the units at work. The share is one input per rig. It is separate from the ship's availability $\epsilon$, from the wind and sea limits, and from launch handling and canopy wear, which are charged only while the kite flies. Of several VTOL kites those still working fly on. Where a kite flies over a row of rotors, the kite and each rotor fail on their own. The fuel store also holds the largest leg with a unit down, without its reserve. A stowed soft kite gives no force, and wings feathered at no lift and stopped rotors stand as drag. Capital, rig mass and maintenance stay charged in both states. No repair cost per event is added, as the maintenance fraction already covers repair and no source gives such a cost (\cref{sec:disc_reliability}).
